\section{Fourier sums and the endpoint summation}\label{sec:fourier}

We now return to periodic functions. In this section the circle has
period one and measure $dy$, as in \Cref{sec:preliminaries}. The Fourier
coefficients use $\chi_k(y)=e^{2\pi iky}$; changing coordinates at the end
recovers precisely the normalization of \Cref{thm:main}.

\subsection{A good-set integral bound}

For a nonnegative periodic function $\rho$, let
\[
 M\rho(x)=\sup_{\substack{I\ni x\\0<|I|\le1}}
                   \frac1{|I|}\int_I\rho(y)\,dy
\]
be its uncentered interval maximal function, with intervals interpreted
on the periodic extension. The elementary interval covering argument
gives
\begin{equation}\label{eq:fourier-maximal-weak}
 \bigl|\{M\rho>\lambda\}\bigr|\le
                 \frac{C}{\lambda}\int_0^1\rho(y)\,dy.
\end{equation}
Indeed, for a compact subset of the level set, choose finitely many
witnessing intervals and a disjoint subfamily whose intervals enlarged
by a fixed factor cover that subset. The sum of the lengths of the
disjoint intervals is at most $\lambda^{-1}\int\rho$, up to an absolute
periodic overlap constant. Inner regularity proves the stated
inequality. In particular, for a mass-one input and $h=\log\log d$,
\[
 G_\rho=\{M\rho\le e^{4h}\},
 \qquad |G_\rho^c|\le Ce^{-4h}=C(\log d)^{-4}.
\]

\begin{theorem}[Truncated integral estimate]\label{thm:good-set}
There are absolute constants $C$ and $d_*$ such that, whenever $d\ge d_*$
and $\rho$ is a nonnegative periodic function with
$\int_0^1\rho=1$ and $\rho\le e^d$ almost everywhere,
\begin{equation}\label{eq:source-18}
       \int_{G_\rho}\sup_{N\ge0}|S_N\rho(x)|\,dx\le Cd.
\end{equation}
\end{theorem}

\begin{proof}
We first pass from \Cref{prop:dyadic} to modulated convolutions, then
identify the periodic Hilbert transform.

\paragraph{The two Haar profiles.}
Extend the following functions by zero outside $[0,1)$:
\[
 \mathfrak h
 =\mathbf1_{[0,1/2)}-\mathbf1_{[1/2,1)},\qquad
 \mathfrak g
 =\mathbf1_{[0,1/4)}-\mathbf1_{[1/4,1/2)}
  -\mathbf1_{[1/2,3/4)}+\mathbf1_{[3/4,1)}.
\]
For a dyadic interval $I=[a,a+r)$ put
$\mathfrak h_I(y)=\mathfrak h((y-a)/r)$ and
$\mathfrak g_I(x)=\mathfrak g((x-a)/r)$.
The shift summand
\[
 \left(\frac1{|I|}\int_I \rho(y)\chi_n(y)\mathfrak h_I(y)\,dy\right)
                \mathfrak g_I(x)
\]
is an admissible summand in \Cref{prop:dyadic}: on the output side choose
$\epsilon_j=r_{j+1}$, so that
$\mathfrak g_I=r_{j+1}r_{j+2}$.

Tile the line by top intervals of length $s\in[1/2,1]$, translate the
tiling by $a\in[0,s)$, and use depths $j=0,\ldots,J-1$ below every top,
where $J\ge1$. Denote the resulting finite-depth shift by
$\mathcal S_{a,s,J}$. This shift need not be periodic, so in the next
estimate $G_\rho$ denotes its representative in $[0,1)$.
At most a fixed number of top intervals meet $[0,1)$.
Each has input average at most $2$, since $\rho$ has mass
one in a period and $s\ge1/2$. At every $x\in G_\rho$, the density cutoff
in \Cref{prop:dyadic} is inactive on all intervals of its spatial path.
Thus the stopped and unstopped shifts agree there, and
\begin{equation}\label{eq:fourier-finite-shift}
 \int_{G_\rho}\sup_{n\in\mathbb Z}
             |\mathcal S_{a,s,J}(\rho\chi_n)(x)|\,dx\le Cd
\end{equation}
uniformly in $a,s,J$.
Here a uniform-probability integral on a top interval is multiplied by
its length to obtain the ordinary spatial integral. Only a bounded
number of top intervals is involved.

\paragraph{Translation and scale averaging.}
Define
\[
 \kappa(u)=\int_{\mathbb R}\mathfrak g(v)\mathfrak h(v-u)\,dv.
\]
For a fixed length $r=s2^{-j}$, uniform translation modulo $s$ is also
uniform translation modulo $r$. Changing variables from an interval's
left endpoint to the output's position $v$ in that interval shows that
the translation average of its full tiling is convolution with
$r^{-1}\kappa(t/r)$. This identity is just Fubini's theorem applied to
the bounded, compactly supported profiles.
For $F=\rho\chi_n$, integrating the translation average in $ds/s$ gives
\[
 \int_{1/2}^1\int_0^s\mathcal S_{a,s,J}F(x)
                 \,\frac{da}{s}\,\frac{ds}{s}=(K_J*F)(x),
\]
where, because the intervals $[2^{-j-1},2^{-j}]$ partition $[2^{-J},1]$,
\begin{equation}\label{eq:fourier-averaged-kernel}
 K_J(t)=\int_{2^{-J}}^1\kappa(t/r)\,\frac{dr}{r^2}.
\end{equation}
The triangle inequality, followed by \Cref{eq:fourier-finite-shift},
therefore yields
\begin{equation}\label{eq:fourier-kernel-bound}
 \int_{G_\rho}\sup_{n\in\mathbb Z}
                  |K_J*(\rho\chi_n)(x)|\,dx\le Cd.
\end{equation}
The measure $ds/s$ on $[1/2,1]$ has mass $\log2$; it has not been
normalized to a probability measure. This fixed factor is absorbed in
$C$. The functions being averaged are finite shift outputs, measurable
in the translation and scale parameters. No measurability of the
auxiliary processing choices used to prove their uniform bound is
required.

\paragraph{The nonzero singular kernel.}
Reflection about $1/2$ fixes $\mathfrak g$ and reverses the sign of
$\mathfrak h$, so $\kappa$ is odd. It is supported in $[-1,1]$ and is
Lipschitz, since it is a finite linear combination of overlaps of
intervals. Moreover
\begin{align}
 c:=\int_0^1\kappa(u)\,du
 &=\int_0^1\mathfrak g(v)
                    \left(\int_0^v\mathfrak h(t)\,dt\right)\,dv
 \notag\\
 &=\frac1{32}-\frac3{32}-\frac3{32}+\frac1{32}
 =-\frac18.                                      \label{eq:fourier-c}
\end{align}
The inner primitive is the tent function $\min(v,1-v)$.
Thus the average in \Cref{eq:fourier-averaged-kernel} has a nonzero
Hilbert-transform component. This calculation agrees with the Haar-shift kernel evaluation in
\cite[Section~3, pp.~1135--1136]{Hytonen2008}.

Put $\delta=2^{-J}$. When $\delta<t\le1$, the substitution $u=t/r$
and the support of $\kappa$ give
\[
 K_J(t)=\frac1t\int_t^1\kappa(u)\,du
       =\frac ct-\frac1t\int_0^t\kappa(u)\,du.
\]
The second term is bounded by an absolute constant; oddness gives
the same conclusion for negative $t$. On $|t|\le\delta$,
\[
 |K_J(t)|\le \|\kappa\|_\infty
                    \int_\delta^1\frac{dr}{r^2}\le C\delta^{-1}.
\]
Compare $K_J$ to
\[
 \pi c\,\cot(\pi t)\,
                \mathbf1_{\{\delta<|t|\le1/2\}}.
\]
On the central annulus the difference is bounded, because
$\pi\cot(\pi t)-1/t$ is bounded on $[-1/2,1/2]$.
On $1/2<|t|\le1$ the original kernel is bounded as well.
Those bounded errors have convolution at most $C$ in absolute value
against any $\rho\chi_n$, using periodicity and mass one.
The inner window contributes at most
\[
 C\delta^{-1}\int_{|t|\le\delta}\rho(x-t)\,dt
 \le C M\rho(x)\le Ce^{4h}\qquad(x\in G_\rho).
\]
It follows from \Cref{eq:fourier-kernel-bound} and $e^{4h}=o(d)$ that
\begin{equation}\label{eq:fourier-truncated-H}
 \int_{G_\rho}\sup_{n\in\mathbb Z}
                      |H_\delta(\rho\chi_n)(x)|\,dx\le Cd,
 \qquad
 H_\delta F(x)=
 \int_{\delta<|t|\le1/2}F(x-t)\cot(\pi t)\,dt.
\end{equation}
Division by the fixed nonzero number $\pi c$ is harmless.

\paragraph{The Hilbert multiplier and Fourier projections.}
For completeness, $H_\delta$ converges in $L^2$ to the periodic
Hilbert transform $H$ whose multiplier is
$-i\,\operatorname{sgn}(k)$ for $k\ne0$ and zero at $k=0$.
To see this directly, oddness reduces its multiplier to a sine
integral. These multipliers are uniformly bounded: write
$\cot(\pi t)=1/(\pi t)+O(1)$, split the sine integral at inverse
frequency, and integrate by parts on the remaining interval.
For a fixed positive integer $k$,
\begin{equation}\label{eq:fourier-symbol}
 \int_{-1/2}^{1/2}\sin(2\pi kt)\cot(\pi t)\,dt=1.
\end{equation}
Indeed, setting $x=\pi t$ and using
\[
 \sin(2kx)=2\sin x\sum_{a=1}^k\cos((2a-1)x)
\]
reduces the integral to the constant term in
$2\cos x\sum_{a=1}^k\cos((2a-1)x)$; all its nonconstant even-frequency
cosines have integral zero on $[-\pi/2,\pi/2]$.
The constant term occurs only for $a=1$ and gives one after division
by $\pi$. Negative frequencies have the opposite sign.
Pointwise multiplier convergence and the uniform bound now imply
$L^2$ convergence by Parseval's identity.

Every $\rho\chi_n$ is bounded and hence belongs to $L^2$.
For a fixed finite set of integers, the maximum of the differences
$|H_\delta(\rho\chi_n)-H(\rho\chi_n)|$ tends to zero in $L^1$,
by the finite sum of their $L^2$ norms. Pass to the limit in
\Cref{eq:fourier-truncated-H} for that finite set and then increase the
set to $\mathbb Z$. Monotone convergence gives
\begin{equation}\label{eq:fourier-H-max}
 \int_{G_\rho}\sup_{n\in\mathbb Z}|H(\rho\chi_n)(x)|\,dx\le Cd.
\end{equation}

Let $P_0$ be projection onto constants and let $\Pi_{\ge0}$ be projection
onto the nonnegative Fourier frequencies. As $L^2$ multiplier identities,
\[
 \Pi_{\ge0}=\frac{I+iH+P_0}{2},
 \qquad
 S_N\rho=
 \chi_{-N}\Pi_{\ge0}(\rho\chi_N)
 -\chi_{N+1}\Pi_{\ge0}(\rho\chi_{-(N+1)}).
\]
They hold almost everywhere simultaneously for all the indicated
integer indices, after removing one null set.
Since $|\widehat\rho(k)|\le1$, their identity and constant terms give
the pointwise upper bound
\[
 \sup_{N\ge0}|S_N\rho|
 \le \rho+1+\sup_{n\in\mathbb Z}|H(\rho\chi_n)|.
\]
Integrating and applying \Cref{eq:fourier-H-max} proves
\Cref{eq:source-18}.
\end{proof}

\subsection{Bounded inputs}

\begin{proposition}\label{prop:fourier-bounded}
Every bounded complex-valued periodic function has almost-everywhere
convergence of its symmetric Fourier partial sums along the full
sequence.
\end{proposition}

\begin{proof}
Let $\|f\|_\infty\le B$. The case $B=0$ is immediate.
Take the Fej\'er means $P_\ell$ of $f$. They are trigonometric
polynomials satisfying $\|P_\ell\|_\infty\le B$ and
$\|f-P_\ell\|_1\to0$. The latter fact follows from the approximate-identity
property of the nonnegative, mass-one Fej\'er kernels and continuity of
translations in $L^1$.

Split $f-P_\ell$ into the positive and negative parts of its real and
imaginary components. Each nonnegative part $q_\ell$ is at most $2B$
and has mass $m_\ell\to0$. If $m_\ell=0$ its partial sums vanish.
Otherwise normalize to $\rho_\ell=q_\ell/m_\ell$ and choose
\[
 d_\ell=d_*+\max\{0,\log(2B/m_\ell)\}.
\]
As $m_\ell\to0$ through positive values, $d_\ell\to\infty$,
$m_\ell d_\ell\to0$, and $\rho_\ell\le e^{d_\ell}$.
\Cref{thm:good-set} and Markov's inequality show, for each $\eta>0$,
\[
 \left|\left\{\sup_N|S_Nq_\ell|>\eta\right\}\right|
 \le C(\log d_\ell)^{-4}
       +\frac{Cm_\ell d_\ell}{\eta}\longrightarrow0.
\]
Applying this to the four component parts proves
$\sup_N|S_N(f-P_\ell)|\to0$ in measure.

For each $\ell$, the partial sums of $P_\ell$ are eventually equal to
$P_\ell$ itself. Consequently the measurable function
\[
 A(x)=\limsup_{N\to\infty}|S_Nf(x)-f(x)|
\]
satisfies
\[
 A(x)\le \sup_N|S_N(f-P_\ell)(x)|+|f(x)-P_\ell(x)|
\]
almost everywhere. Both terms on the right tend to zero in measure.
For each $\eta>0$ this bounds $|\{A>\eta\}|$ by a quantity tending to
zero. Thus $A=0$ almost everywhere, as required.
\end{proof}

\subsection{Summation of the height layers}

\begin{proof}[Proof of \Cref{thm:main}]
It suffices to work in unit-period coordinates. Write
\[
 f=f_{\mathrm{bd}}+\sum_{j\ge1}f_j,\qquad
 f_{\mathrm{bd}}=f\,\mathbf1_{\{|f|<e^2\}},\qquad
 f_j=f\,\mathbf1_{\{e^{2^j}\le|f|<e^{2^{j+1}}\}}.
\]
For each $j$, split $f_j$ into four nonnegative parts with coefficients
$1,-1,i,-i$, using the positive and negative parts of its real and
imaginary components. Consider one such sequence $q_j$, and write
$m_j=\int q_j$. Discard its zero-mass members. Put
\[
 \rho_j=q_j/m_j,\qquad
 d_j=2^{j+1}+\log^-m_j.
\]
Then $\int\rho_j=1$ and $\rho_j\le e^{d_j}$.
For all sufficiently large $j$, \Cref{thm:good-set} applies.

We claim
\begin{equation}\label{eq:fourier-layer-cost}
                  \sum_j m_jd_j<\infty.
\end{equation}
On the support of $q_j$, $\log|f|\ge2^j$, and $q_j\le|f|$.
Disjointness of the layers therefore gives
\[
 \sum_j2^{j+1}m_j
 \le2\int_{\{|f|\ge e^2\}}|f|\log|f|<\infty.
\]
For the additional entropy term, if $m_j\ge e^{-2^j}$, then
$m_j\log^-m_j\le2^jm_j$. If $m_j<e^{-2^j}$, the monotonicity of
$t\log(1/t)$ on $(0,e^{-1})$ gives
$m_j\log(1/m_j)\le2^je^{-2^j}$. Both bounds are summable.
This proves \Cref{eq:fourier-layer-cost}.

Let $G_j$ be the good set associated to $\rho_j$. Since
$\log d_j\ge(j+1)\log2$,
\[
 \sum_j|G_j^c|\le C\sum_j(\log d_j)^{-4}<\infty.
\]
Borel--Cantelli implies that almost every point belongs to all but
finitely many $G_j$. On the other hand, Tonelli's theorem and
\Cref{eq:source-18,eq:fourier-layer-cost} give
\[
 \int\sum_j\mathbf1_{G_j}(x)\sup_N|S_Nq_j(x)|\,dx
 \le C\sum_jm_jd_j<\infty,
\]
where a finite number of initial layers may be omitted.
Therefore, at almost every point, the partial-sum maxima of the tail
layers are absolutely summable.
Every individual $q_j$ is bounded, so \Cref{prop:fourier-bounded}
also gives its full-sequence limit and the finiteness of its
partial-sum maximum outside a null set. Restoring finitely many
omitted layers causes no difficulty.
Intersect these full-measure sets for the four component sequences
and for $f_{\mathrm{bd}}$.

The layers are summable in $L^1$. For each fixed $N$, Fourier
coefficients can consequently be summed term by term; equivalently,
the finite-rank operator $S_N$ is bounded from $L^1$ to $L^\infty$.
Thus
\[
 S_Nf=S_Nf_{\mathrm{bd}}+
       \sum_{j\ge1}S_Nf_j
\]
with the four component parts understood in the last sum.
At every point of the full-measure set just constructed, the absolute
sum of their partial-sum maxima is finite. Its tail tends to zero
uniformly in $N$. Take the limit first on a finite collection of
layers, using \Cref{prop:fourier-bounded}, and then let the collection
increase. The result is
$S_Nf(x)\to f_{\mathrm{bd}}(x)+\sum_jf_j(x)=f(x)$.
All excluded sets are measurable, and only countably many are used.
Changing back to $x=2\pi y$ proves the stated theorem with normalized
measure $dx/(2\pi)$.
\end{proof}

\subsection{The maximal consequence}

\begin{corollary}[Weak-$L^1$ maximal bound]\label{cor:fourier-weak-l1}
On $\mathbb T=\mathbb R/(2\pi\mathbb Z)$ with $d\mu=dx/(2\pi)$, set
$\Phi(t)=t\log(2+t)$ and define the Luxemburg norm by
\[
 \|f\|_\Phi
 =\inf\left\{b>0:\int_{\mathbb T}\Phi(|f|/b)\,d\mu\le1\right\}.
\]
For the ordinary symmetric partial sums of \Cref{thm:main}, put
$S^*f=\sup_{N\ge0}|S_Nf|$. There is an absolute constant $C$ such that
\begin{equation}\label{eq:fourier-weak-l1}
 \sup_{\lambda>0}\lambda\,\mu\{S^*f>\lambda\}
 \le C\|f\|_\Phi
 \qquad(f\in L\log L(\mathbb T;\mathbb C)).
\end{equation}
\end{corollary}

\begin{proof}
The real Luxemburg space $B=L^\Phi(\mathbb T;\mathbb R)$ is Banach.
Since $\Phi(2t)\le4\Phi(t)$, its underlying set is exactly the real
$L\log L$ class used in \Cref{thm:main}; moreover
$\|f\|_1\le\|f\|_\Phi/\log2$.
Translations are isometries of $B$ by invariance of $\mu$.
Each $S_N$ is a real linear operator commuting with translations and
is bounded from $B$ into itself: indeed,
$\|S_Nf\|_\infty\le(2N+1)\|f\|_1$, and $L^\infty$ embeds continuously
in $B$. These bounds need not be uniform in $N$.
By \Cref{thm:main}, $S^*f$ is finite almost everywhere for every $f\in B$.
Stein's theorem \cite[Appendix, Theorem~10]{Stein1961}, applied to the
sequence $S_0,-S_0,S_1,-S_1,\ldots$, therefore gives
\Cref{eq:fourier-weak-l1} for real inputs; the signed supremum in that
theorem is precisely $S^*f$ for this sequence.
For $f=u+iv$, use $S^*f\le S^*u+S^*v$ and
$\|u\|_\Phi,\|v\|_\Phi\le\|f\|_\Phi$, splitting the level at
$\lambda/2$. This proves the complex-valued assertion after enlarging
$C$ by a fixed factor.
\end{proof}
